% ============================================================================
\chapter{B-Splines and Learnable Activation Functions}
\label{ch:splines}
% ============================================================================

This chapter introduces the mathematical foundation for the third
component of the hybrid architecture: learnable functions that are
smooth during training but can be discretised into lookup tables for
deployment. The key tool is the B-spline, a flexible curve defined by
a small set of control points.


% ----------------------------------------------------------------------------
\section{Why Learnable Functions?}
\label{sec:why-learnable}
% ----------------------------------------------------------------------------

In a conventional neural network, activation functions (ReLU, sigmoid,
etc.) are fixed by the architect. The \concept{Kolmogorov--Arnold
Network} (KAN)~\cite{liu2024kan} proposes a radical alternative: let
the network \emph{learn} its activation functions during training.

This is motivated by the Kolmogorov--Arnold representation theorem~\cite{kolmogorov1957representation}:

\begin{theorem}[Kolmogorov--Arnold, 1957~\cite{kolmogorov1957representation}]
\label{thm:ka}
Any continuous function $f : [0,1]^n \to \mathbb{R}$ can be written as:
\begin{equation}
  f(x_1, \ldots, x_n) = \sum_{q=0}^{2n} \Phi_q\!\left(
    \sum_{p=1}^{n} \phi_{q,p}(x_p)\right)
  \label{eq:ka-theorem}
\end{equation}
where $\phi_{q,p} : [0,1] \to \mathbb{R}$ and
$\Phi_q : \mathbb{R} \to \mathbb{R}$ are continuous functions.
\end{theorem}

The theorem says that multivariate functions can be decomposed into
compositions and sums of \emph{univariate} functions ($\phi$ and
$\Phi$). If we can learn these univariate functions well, we can
represent any continuous mapping.

\begin{analogy}
Traditional neural networks are like a factory with fixed machines
(ReLU, sigmoid) connected by adjustable conveyor belts (weights). The
KAN approach is like a factory where the machines themselves are
reprogrammable---each machine learns the optimal transformation for
its position in the assembly line.
\end{analogy}


% ----------------------------------------------------------------------------
\section{B-Splines: A Primer}
\label{sec:bsplines}
% ----------------------------------------------------------------------------

A \concept{B-spline} (basis spline) is a smooth, piecewise polynomial
curve defined by \emph{control points} and
\emph{knots}~\cite{deboor1978splines}.

\begin{definition}[B-spline curve]
Given $n+1$ control points $c_0, c_1, \ldots, c_n$ and a knot vector
$\bvec{t} = (t_0, t_1, \ldots, t_{n+k})$, the B-spline of degree
$k{-}1$ is:
\begin{equation}
  S(x) = \sum_{i=0}^{n} c_i \cdot B_{i,k}(x)
  \label{eq:bspline}
\end{equation}
where $B_{i,k}(x)$ are the B-spline basis functions, defined
recursively by the Cox--de~Boor formula.
\end{definition}

For our purposes, the important properties are:

\begin{enumerate}
  \item \textbf{Smooth:} B-splines are continuously differentiable
    (up to degree $k{-}2$). Gradients flow through them.
  \item \textbf{Local control:} Moving one control point affects
    only a nearby portion of the curve. This makes learning stable.
  \item \textbf{Compact representation:} A curve is defined by a
    small number of control points (typically 8--16), regardless of
    how many points you evaluate it at.
  \item \textbf{Evaluable:} Given an input $x$, the output $S(x)$
    can be computed in $O(k)$ operations.
\end{enumerate}

\begin{figure}[htbp]
\centering
\begin{tikzpicture}[>=Stealth]
  \begin{axis}[
    width=10cm, height=5cm,
    xlabel={Input $x$}, ylabel={Output $S(x)$},
    xmin=0, xmax=1, ymin=-0.5, ymax=1.5,
    grid=major, grid style={bitgray!20},
    every axis label/.style={font=\small\sffamily},
    tick label style={font=\small},
  ]
    % A smooth learned function
    \addplot[very thick, bitblue, smooth, domain=0:1, samples=100]
      {0.2 + 0.8*sin(deg(3.14*x)) - 0.3*sin(deg(6.28*x))};

    % Control points
    \addplot[only marks, mark=*, mark size=3pt, bitorange]
      coordinates {
        (0, 0.2) (0.125, 0.5) (0.25, 0.9)
        (0.375, 1.1) (0.5, 0.8) (0.625, 0.4)
        (0.75, 0.1) (0.875, 0.3) (1.0, 0.5)
      };

    \node[font=\scriptsize\sffamily, bitorange] at (axis cs:0.85, 1.2)
      {control points};
    \node[font=\scriptsize\sffamily, bitblue] at (axis cs:0.15, 1.2)
      {learned curve};
  \end{axis}
\end{tikzpicture}
\caption{A B-spline curve defined by 9 control points (orange dots).
  The smooth curve (blue) passes near the control points. During
  training, gradient descent adjusts the control points to minimise
  the loss.}
\label{fig:bspline}
\end{figure}


% ----------------------------------------------------------------------------
\section{KAN: Putting Splines on Edges}
\label{sec:kan}
% ----------------------------------------------------------------------------

In a traditional neural network, computation is:
\begin{equation}
  y = \sigma(w \cdot x + b)
  \label{eq:trad-neuron-recap}
\end{equation}
The \emph{weight} $w$ on the edge is a single number (linear
scaling), and the \emph{activation} $\sigma$ on the node is a fixed
function.

In a KAN, this is reversed:
\begin{equation}
  y = \sum_{i} \phi_i(x_i)
  \label{eq:kan-neuron}
\end{equation}
The \emph{edge} carries a learnable function $\phi_i$ (a B-spline),
and the \emph{node} is a simple summation.

\begin{figure}[htbp]
\centering
\begin{tikzpicture}[
  input/.style={circle, draw, minimum size=0.6cm, font=\small},
  neuron/.style={circle, draw, minimum size=0.9cm, font=\small,
    fill=bitblue!10},
  spline/.style={draw, rounded corners, fill=bitorange!10,
    minimum width=1.2cm, minimum height=0.5cm, font=\tiny\sffamily},
  >=Stealth
]
  % Traditional
  \node[font=\small\sffamily\bfseries] at (0, 3) {Traditional};
  \node[input] (t1) at (-1, 1.5) {$x_1$};
  \node[input] (t2) at (-1, 0) {$x_2$};
  \node[neuron] (tn) at (1.5, 0.75) {$\sigma(\Sigma)$};
  \draw[->, thick] (t1) -- (tn) node[midway, above, font=\tiny, sloped] {$w_1$};
  \draw[->, thick] (t2) -- (tn) node[midway, below, font=\tiny, sloped] {$w_2$};

  % KAN
  \node[font=\small\sffamily\bfseries] at (6, 3) {KAN};
  \node[input] (k1) at (5, 1.5) {$x_1$};
  \node[input] (k2) at (5, 0) {$x_2$};
  \node[neuron] (kn) at (8.5, 0.75) {$\Sigma$};
  \node[spline] (s1) at (6.7, 1.5) {$\phi_1$};
  \node[spline] (s2) at (6.7, 0) {$\phi_2$};
  \draw[->, thick] (k1) -- (s1);
  \draw[->, thick] (s1) -- (kn);
  \draw[->, thick] (k2) -- (s2);
  \draw[->, thick] (s2) -- (kn);
\end{tikzpicture}
\caption{Traditional neuron vs.\ KAN neuron. In a traditional neuron,
  edges carry scalar weights and the node applies a fixed activation.
  In a KAN, edges carry learnable functions (B-splines) and the node
  simply sums.}
\label{fig:kan-vs-trad}
\end{figure}

Each learnable function $\phi_i$ is a B-spline~\cite{deboor1978splines} with (say) 8 control
points. The control points are the trainable parameters. During
training, backpropagation adjusts these control points to minimise
the loss.


% ----------------------------------------------------------------------------
\section{Relevance to Our Architecture}
\label{sec:relevance}
% ----------------------------------------------------------------------------

Our hybrid architecture does not use a full KAN. Instead, it borrows
the key idea: \emph{learnable univariate functions that can be
compiled to lookup tables.}

The specific application is in the scoring layer. After the LGN
produces binary outputs, we need to aggregate them into class scores.
A simple approach is popcount (count the 1s for each class group),
but a learned scoring function can be more effective:

\begin{equation}
  s_k = \phi_k\!\left(\sum_{j \in \text{group}_k} g_j\right)
  \label{eq:scoring}
\end{equation}

Here, $\sum g_j$ is the popcount of the gate outputs for class $k$,
and $\phi_k$ is a learned function that maps this count to a score.
During training, $\phi_k$ is a B-spline. After training, it is
compiled into a lookup table.

\begin{keyinsight}
The computation-memory trade-off at work: instead of computing a
complex nonlinear function at inference time, we precompute it for
all possible input values and store the results in a table. The
function's complexity (how many control points, how many spline
segments) affects only the \emph{training} cost. The inference cost
is always a single table lookup, regardless of the function's
complexity.
\end{keyinsight}


% ----------------------------------------------------------------------------
\section{Training the Scoring Functions}
\label{sec:scoring-training}
% ----------------------------------------------------------------------------

The scoring functions can be trained jointly with the LGN (during
Phase~2) or in a separate Phase~2b after the LGN is discretised.
The latter approach is simpler:

\begin{enumerate}
  \item Discretise the LGN (fix all gate types).
  \item For each training image, compute the popcount vector
    $(n_0, n_1, \ldots, n_9)$ where $n_k$ is the number of output
    gates for class $k$ that fire.
  \item Train the scoring functions $\phi_0, \ldots, \phi_9$ on
    these popcount vectors using cross-entropy loss and gradient
    descent.
  \item The only trainable parameters are the B-spline control points
    (8--16 per function, 80--160 total for 10 classes).
\end{enumerate}

This is a tiny optimisation problem that converges in a few epochs.


% ----------------------------------------------------------------------------
\section{Why B-Splines and Not Polynomials?}
\label{sec:why-splines}
% ----------------------------------------------------------------------------

One might ask: why not use a simple polynomial $\phi(x) = a_0 +
a_1 x + a_2 x^2 + \cdots$?

\begin{enumerate}
  \item \textbf{Numerical stability:} High-degree polynomials
    oscillate wildly (Runge's phenomenon). B-splines are locally
    defined and stable.
  \item \textbf{Local control:} Adjusting one control point of a
    B-spline affects only a nearby region. Adjusting one polynomial
    coefficient affects the entire curve.
  \item \textbf{Discretisation quality:} A B-spline with 8 control
    points can be sampled at, say, 64 points to create a lookup table
    with very small approximation error. A polynomial of degree 8
    might require more samples to capture its oscillations.
\end{enumerate}

\begin{engineernote}
For the scoring functions in our architecture, the input is a
popcount value in the range $[0, n_{\text{gates}}]$. With $n_{\text{gates}} = 200$
per class, the lookup table has 201 entries. Each entry is an 8-bit
or 16-bit integer. The total table size is $10 \times 201 \times 2 =
4{,}020$ bytes---negligible.
\end{engineernote}


% ----------------------------------------------------------------------------
\section{Summary}
\label{sec:ch10-summary}
% ----------------------------------------------------------------------------

This chapter introduced B-splines as learnable functions that bridge
the gap between continuous training and discrete deployment:

\begin{enumerate}
  \item The Kolmogorov--Arnold theorem~\cite{kolmogorov1957representation} guarantees that compositions
    of univariate functions can represent any continuous mapping.
  \item B-splines~\cite{deboor1978splines} provide smooth, locally controllable, compactly
    parameterised univariate functions.
  \item During training, B-splines support standard backpropagation
    through their control points.
  \item After training, each B-spline can be sampled and stored as a
    lookup table---converting a continuous function into a discrete
    memory access.
\end{enumerate}

The next chapter shows how to compile these functions into lookup
tables and analyses the accuracy-efficiency trade-off of
discretisation.
